\subsection{主要贡献总结}
\label{sec:conclusion-contributions}

本文对激光超声无损检测做出以下四项主要贡献：

\begin{enumerate}
\item \textbf{PMDF-Net多域融合架构。} 一种双分支编码器-解码器，联合处理时域波形与时频表示，提取互补特征以增强信号去噪保真度。

\item \textbf{物理约束损失函数。} 一种保特征损失 formulation，强制执行到达时间、振幅与波形形态一致性，确保去噪信号保留用于缺陷表征的诊断信息。

\item \textbf{减平均训练协议。} 一种配对训练策略，使用低平均次数输入 ($N=16$) 对抗高平均次数目标 ($M=256$)，使实际部署无需在检测过程中进行大量信号平均。

\item \textbf{跨域泛化能力。} 展示了从完整铝试样到含缺陷工件的跨域迁移，以及跨厚度变化的泛化，验证了该方法在训练条件之外的真实检测场景中的适用性。
\end{enumerate}

这些贡献确立了一个高效激光超声检测框架，在减少平均次数需求的同时，保留并超越了传统平均方法在缺陷检测任务中的性能。
